
Da wir diese Additivität aus der proximity structure herleiten wollen, müssen wir Eigenschaften definieren, aus der man die Additivität herleiten kann. Dazu benötigen wir Folgenden Begriff:

\begin{defi}
Für $a, b, c \in  A$ gilt die dreistellige Relation $(a, b, c \rangle$ wenn für alle $a', b', c' \in A$ gilt:
\begin{enumerate}
  \item wenn $(a ,b) \succsim (a', b')$ und  $(a', c') \succsim (a, c)$, dann gilt  $(b', c') \succsim (b, c)$
  \item Wenn beide der Ungleichungen strikt sind, ist die resultierende Ungleichung auch strikt
\end{enumerate}
Weiter gilt $\langle abc \rangle$ wenn $(a, b, c \rangle$ und $(c, b, a\rangle$ beide gelten.
\end{defi}


Die Bedingungen, die notwendig sind, um eine Darstellung durch eine Metrik mit additiven Segmenten zu ermöglichen, sind in der folgenden Definition und dem Satz danach zusammengefasst. Der zweite Teil der Definition wird erst später, für das Hauptresultat (Satz \ref{lp-eindeutig}) relevant, im Prinzip geht es darum dass man sich Grenzwerte und Cauchy Folgen anschauen kann, also der normale Begriff von Vollständigkeit.

\begin{defi}[segmentally additive proximity structure]~\\
Eine proximity structure $\langle A, \succsim \rangle$, mit einer dreifachen Relation $\langle \rangle$, die in Bezug auf $\succsim$ definiert ist, ist segmentally additive genau dann, wenn die folgenden zwei Axiome für alle $a, c, e, f \in A$ gelten:

\begin{enumerate}
    \item Wenn $(a, c) \succsim (e, f)$, dann gibt es ein $b \in A$, so dass $(a, b) \sim (e, f)$ und $\langle abc \rangle$. (segmental solvability)
    \item Wenn $e \neq f$, dann gibt es $b_0, \ldots, b_n \in A$, so dass $b_0 = a, b_n = c$ und $(e, f) \succsim (b_{i-1}, b_i), i = 1, \ldots, n$. (connectedness)
\end{enumerate}

Eine segmentally additive proximity structure ist genau dann vollständig, wenn die folgenden zwei Axiome für alle $a, c, e, f \in A$ gelten:

\begin{enumerate}
    \item[3.] Wenn $e \neq f$, dann gibt es $b, d \in A$, wobei $b \neq d$, so dass $(e, f) \succ (b, d)$. (Nichtdiskretheit der Abstände)
    \item[4.] Sei $a_1, \ldots, a_n, \ldots$ eine Folge von Elementen von $A$, so dass für alle $e, f \in A$, wenn $e \neq f$, dann $(e, f) \succsim (a_i, a_j)$ für alle bis auf endlich viele Paare $(i, j)$. Dann gibt es ein $b \in A$, so dass für alle $e, f \in A$, wenn $e \neq f$, dann $(e, f) \succsim (a_i, b)$ für alle bis auf endlich viele $i$. (Grenzwerte für Cauchy Folgen)
\end{enumerate}

\end{defi}


\addcontentsline{toc}{subsection}{seg add representability}
